\begin{theorem}[The Hodge obstruction to normal-crossing limits]\hypertarget{tag-030Z}{}\label{push:src:B.4.3}
\StatementRecord{push:src:B.4.3}{proved}{030Z}{\SeriesRef{push:src:B.4.33}; \citet[Theorem 1.4, preprint p.~7, and Example 2.1, p.~13]{Nishiguchi1988IIRoot}; \citet[Theorem 7.1]{IstratiEtAl2022Root}}

Let $g:\mathcal Z\to\Delta_s$ be a proper holomorphic map from a
smooth complex fourfold, smooth over $\Delta_s^*$, with reduced
normal-crossings central fibre $D$. Suppose its nearby fibres are
biholomorphic to the six-sphere members of the specified family,
after any finite base change $w=s^N$. Then
\[
 H^1(D;\C)=0,\qquad h^1(D,\OO_D)\geq1.
\]
Consequently the natural map
$H^1(D;\C)\to H^1(D,\OO_D)$ is not surjective. In particular,
$D$ cannot be the analytification of a proper complex algebraic
variety with normal-crossings singularities.

If $D$ has \emph{simple} normal crossings, at least one of its
nonempty closed strata fails the $\partial\bar\partial$-lemma.
In particular, no such simple normal-crossings filling can have
all irreducible components K\"ahler or all irreducible components
Moishezon. These obstructions hold for every $N\geq1$.
\end{theorem}

\begin{proof}
Theorem~\ref{push:src:B.4.2} gives a surjection from the fundamental group of a nearby fibre onto $\pi_1(D)$. Since that fibre is diffeomorphic to $S^6$, it follows that $\pi_1(D)=1$ and $H^1(D;\mathbb C)=0$.

For a compact complex threefold $X$ diffeomorphic to $S^6$, the groups $H^2(X;\mathbb Z)$ and $H^4(X;\mathbb Z)$ vanish, so $c_1(X)=c_2(X)=0$. The analytic Riemann--Roch formula, obtained from \citet[Theorem~4.3, pp.~564--565]{AtiyahSinger1968IIIE}, gives
\[
 \chi(X,\mathcal O_X)=\frac1{24}\langle c_1(X)c_2(X),[X]\rangle=0.
\]
Every holomorphic three-form $\Omega$ is closed and is exact because $H^3(X;\mathbb C)=0$. Stokes' theorem gives $\int_X i\Omega\wedge\overline\Omega=0$. The integrand is pointwise nonnegative for the complex orientation, so $\Omega=0$. Serre duality therefore gives $h^3(X,\mathcal O_X)=0$, and connectedness gives $h^0(X,\mathcal O_X)=1$. Hence
\[
 h^1(X,\mathcal O_X)=1+h^2(X,\mathcal O_X)\geq1.
\]
Locally along the reduced normal-crossing divisor, $g$ has the form $s=z_1\cdots z_k$ after absorbing a holomorphic unit. This makes $\mathcal O_{\mathcal Z}$ torsion-free over the local discrete valuation ring of the disc and hence flat over it. Upper semicontinuity for the proper flat map gives $h^1(D,\mathcal O_D)\geq1$; see \citet[\S7, Satz~5, pp.~63--64]{Grauert60}.

For a proper complex algebraic variety with Du Bois singularities, the canonical map from singular cohomology to coherent cohomology is surjective in every degree; see \citet[\S1.D, p.~120]{GrafKovacs2014Root}. Normal-crossing singularities are Du Bois, including self-intersections of a global component. Thus algebraicity and properness of $D$ would make $H^1(D;\mathbb C)\to H^1(D,\mathcal O_D)$ surjective, contradicting the two groups just computed.

Now suppose $D$ is simple normal crossing and each nonempty closed stratum satisfies the $\partial\bar\partial$-lemma. These are exactly the strata in condition (K1) of \citet[Definition~2.3, p.~9]{Chen2026HodgeRoot}. For a proper semistable degeneration with smooth total space and this condition, \citet[Theorem~2.4(4), p.~10]{Chen2026HodgeRoot} gives $E_1$ degeneration of the Hodge--de Rham spectral sequence on a nearby fibre. In total degree one this yields
\[
 0=b_1(X_s)=h^0(X_s,\Omega^1_{X_s})+h^1(X_s,\mathcal O_{X_s})\geq1,
\]
a contradiction.

A closed smooth submanifold of a compact K\"ahler manifold is K\"ahler. A closed stratum of a Moishezon component is also Moishezon: choose a projective modification, take a component of the inverse image dominating the stratum, and intersect with general relative hyperplanes until the induced proper map is generically finite. This gives a generically finite projective cover and hence maximal algebraic dimension for the stratum. Smooth Moishezon manifolds satisfy the $\partial\bar\partial$-lemma, as recalled in \citet[\S2.1, p.~8]{Chen2026HodgeRoot}. Either stated hypothesis on the irreducible components therefore gives condition (K1), proving the last assertion. A finite parameter cover preserves the biholomorphism type of every nonzero fibre, so the argument applies for every $N\geq1$.
\end{proof}
